\documentclass[10pt,a4paper]{amsart}
\usepackage[margin=19mm]{geometry}
\usepackage{amsmath,amssymb,mathtools}
\usepackage[scr=rsfs,cal=euler]{mathalfa}
\usepackage{xfrac}
\usepackage[unicode,hidelinks,bookmarks=false]{hyperref}
\newcommand{\ie}{i.e.\ }
\newcommand{\cf}{cf.\ }
\newcommand{\resp}{respectively\ }
\newcommand{\Subsection}{\subsection}
\providecommand{\nobreakdash}{\nobreak\hbox{-}}
\setlength{\emergencystretch}{2em}
\multlinegap0pt
\usepackage{bm}
\usepackage{bbm}
\usepackage{dsfont}
\usepackage{xspace}
\usepackage{cancel}
\usepackage{mathtools}
\usepackage{amsthm}
\makeatletter
\@ifundefined{theorem}{\newtheorem{theorem}{Theorem}}{}
\@ifundefined{lemma}{\newtheorem{lemma}[theorem]{Lemma}}{}
\makeatother
\def\R{\mathbb{R}}
\newcommand{\cF}{\mathcal F}
\title[Wave packet systems and connections to spectral analysis of ]{Wave packet systems and connections to spectral analysis of limiting operators}
\author{Kevin Hughes, Arie Israel, Azita Mayeli}
\date{Source: arXiv:2601.06945v1 (CC BY 4.0)}
\begin{document}
\maketitle

\noindent\emph{Source and licence.} Wave packet systems and connections to spectral analysis of limiting operators. Version arXiv:2601.06945v1 (CC BY 4.0), distributed under the Creative Commons Attribution 4.0 International licence, as recorded in the arXiv licence metadata for this exact version and shown on the abstract page https://arxiv.org/abs/2601.06945v1. Full rights notice: licenses/arxiv-2601.06945-CC-BY-4.0.txt.\par\smallskip

\noindent\textit{Editorial scope.} Counted unit: the Lemma whose source label is \texttt{packing\_spectral:lem} in the article (source0.tex lines 486--488, source label \texttt{packing\_spectral:lem}), delivered together with the article's own complete original proof of it (source0.tex lines 489--518, one \texttt{proof} environment of 30 lines). No mathematics is written, repaired or re-derived: statement and proof are the article's own text line for line. Required context retained uncounted: none. Every definition, display and label of those lines is retained unchanged. Omitted from this package and not redistributed: 1--485, 519--683. Nothing inside a retained excerpt is omitted or altered: every retained body is delivered exactly as the article's lines are written, so no omission receipt of any kind accompanies this package. Citations: the retained excerpts carry no citation command at all, so no bibliography entry of the article is transcribed and no reference is redistributed. Reference adaptation: none was needed and none was made; every reference inside the delivered unit resolves inside it, so no unresolved reference or citation is produced. Printed slips kept literal: no printed word, symbol, hypothesis or number of the counted statement or of its proof was altered. Layout and class: the article's own class is \texttt{IEEEtran} with packages including amsmath, amssymb, amsfonts, algorithmic, graphicx, textcomp, xcolor, biblatex, amsthm, latexsym, epsfig, enumerate, multicol, wasysym; the wrapper is the lane's pinned \texttt{amsart} editorial wrapper at 10pt on A4, which restates the theorem-like environments the retained text needs and adds the article's own macro definitions that the retained text needs (\texttt{\textbackslash R}, \texttt{\textbackslash cF}), with extra wrapper packages (bm, bbm, dsfont, xspace, cancel, mathtools). The delivered numbering is the unit's own. Assets: no retained span contains a figure environment, a graphics-inclusion command, a PDF-page inclusion, a table or a TikZ picture; no image file is copied into this package, the package declares no asset dependency and no image file is redistributed with it. Licence: version arXiv:2601.06945v1 of this article is distributed under the Creative Commons Attribution 4.0 International licence, as recorded in the arXiv licence metadata for this exact version and shown on the abstract page https://arxiv.org/abs/2601.06945v1. A full rights notice is delivered at licenses/arxiv-2601.06945-CC-BY-4.0.txt. Characters: every retained excerpt is pure ASCII, so the delivered engine, XeLaTeX with its default Latin Modern text font, reports no missing character.
\begin{lemma}\label{packing_spectral:lem}
Let $\cF = \{\psi_\nu\}$ be an $\epsilon$-packing  for $F \times S$, with $n := \#\cF$, and $\epsilon < \frac{1}{2n}$. Then $\lambda_{n}(P_F B_S P_F) > 1 - 5 \epsilon n^{1/2}$.
\end{lemma}

\begin{proof}
We write $\cF = \{\psi_1,\cdots,\psi_n\}$, and let $V \subset L^2(\R^d)$ be the subspace spanned by $\cF$. We claim that $\cF$ is a linearly independent set. Indeed, let $G$ be the $n \times n$ Gram matrix of $\cF$, i.e., $G_{\nu \nu'} = \langle \psi_\nu,\psi_{\nu'}\rangle$, $1 \leq \nu,\nu' \leq n$. The diagonal entries of $G$ are $1$, and the off-diagonal entries of $G$ have magnitude at most $\epsilon$. Let $I$ denote the $n \times n$ identity matrix. Observe that the Frobenius norm of $I - G$ is at most $\epsilon n < 1/2$. The spectral norm is bounded by the Frobenius norm, and so, $\| I - G \| < 1/2$, hence $G$ is invertible, and so $\cF$ is linearly independent. Thus, $\dim V = n$. 

By the variational characterization of eigenvalues (the max-min theorem), the bound $\lambda_n(F,S) > 1 - C \epsilon n^{1/2}$ is implied by
% \begin{align}
%  \|P_F B_S P_F \psi\|  > (1- C \epsilon n) \| \psi \| \quad \forall \psi \in V,  ~ \psi \neq 0.
% \end{align}
\[
\frac{\|P_F B_S P_F \psi\|}{\| \psi \|} > 1- C \epsilon n^{1/2} \quad \forall \psi \in V,  ~ \psi \neq 0.
\]

The claim that $\{\psi_\nu\}$ is $\epsilon$-concentrated on $F \times S$ implies that $\|(Id - P_F) \psi_\nu \| \leq \epsilon$ and $\|(Id - B_S) \psi_\nu \| \leq \epsilon$ for all $\nu$. Observe that $Id - P_F B_S P_F = (Id - P_F) + P_F(Id - B_S) + P_F B_S (Id - P_F)$. Since $P_F$ and $B_S$ are projections, by the triangle inequality we obtain $\|(Id - P_F B_S P_F) \psi_\nu \| \leq 3 \epsilon$.

Let $\psi \in V$, and write $\psi = \sum_\nu a_\nu \psi_\nu$. By the triangle inequality and the estimation in the previous paragraph, we have $\| (Id - P_FB_SP_F) \psi \| \leq 3 \epsilon \sum_{\nu} |a_\nu|$, which implies that
\begin{align*}
\| P_F B_S P_F \psi \| \geq \| \psi \| - 3 \epsilon \sum_\nu |a_\nu| \geq \| \psi\| - 3 \epsilon n^{1/2} 
 % \left(\sum_\nu |a_\nu|^2\right)^{1/2}.
 \sqrt{\sum_\nu |a_\nu|^2}.
\end{align*}
 
By expanding out the inner product, we write
\[
\| \psi \|^2 = \sum_{\nu} |a_\nu|^2  + \sum_{\nu\neq \mu} a_\nu \overline{a_\mu} \langle \psi_\nu, \psi_\mu \rangle \geq (1 - n\epsilon) \sum_\nu |a_\nu|^2,
\]
where in the last inequality we used $\left|\sum_{\nu \neq \mu} a_\nu \overline{a_\mu} \right| \leq (\sum |a_\nu| )^2 \leq n \sum_\nu |a_\nu|^2$, as well as the approximate orthogonality of $\cF$. Combining the previous estimates,
\[
\frac{\| P_F B_S P_F \psi \|}{\| \psi \|}  \geq 1 - 3 \epsilon n^{1/2}(1-n\epsilon)^{-1/2} \geq 1 - 5 \epsilon n^{1/2}.
\]
This estimate holds for any nonzero $\psi \in V$. This concludes the proof of the lemma.
\end{proof}

\end{document}
